\section{Information, prior results and retained conclusions}\label{sec:comparison}
\subsection{What has changed in the experiment}
The previous local-value theorem used active local value and distance
access: certified Cartesian chart coverage, boundary-height enclosures
at arbitrarily prescribed tolerance, and enclosures for the distance to
\emph{all} solid near a candidate flight. Its effectiveness was relative
to complex-analytic and geometric numerical priors. Its explicit uniform
short-arc fingerprint could be extremely large. Those statements are
retained with their original assumptions; a symbolic finite expression
was not a practical sensor design.

The present experiment requests none of those geometric values. It
uses freely positioned short launches and localized collision coordinates,
and derives confidence geometry from a finite random cloud. Its local
spatial neighborhood is large enough to sample the complete encountered
bodies. Thus this is a change in access and range, not a claim that the
new data are a subset of the old short-patch data. The number of launches
and the required localization tolerance are both displayed. Independent
resettable launches remain stronger than observing one passive trajectory.
The fixed separation and genericity bounds remain priors, not estimates.

The gain in regularity is correspondingly precise. Whole boundary images
are now sampled, so the certificate works with finite $C^{6,\beta}$ bounds
and needs no holomorphic strip, graph disk or continuation modulus. The
class is not merely analytic bodies described with different notation.
For example, at any member with strict margins, sufficiently small
$C^{6,\beta}$ perturbations of its support functions preserve all the
finitely many bounded-copy inequalities and the shape and symmetry gaps.
One may choose a smooth periodic bump, flat at the edge of its support
and nonzero inside, which is not real analytic. These perturbations give
nonanalytic physical members wherever the original strict-margin class
is nonempty. The earlier registered symmetric-table and analytic
short-contact results are not restricted by this new acquisition theorem.

\subsection{Nearest inverse-geometry comparisons}
The canonical relation generated by $W$ has momenta $-W_s,W_t$.
Knowing that relation gives $W$ only up to a constant on its local branch;
an absolute action value and the normalization $A$ are additionally
needed to recover the two densities. Conversely \eqref{eq:action-area}
recovers these from the densities. Differentiation of travel-time data
is a classical inverse-geometry mechanism, not a new abstract principle
claimed here.

Stefanov--Uhlmann--Vasy~\cite{SUV} determine a metric near a strictly
convex boundary point from local boundary distance, modulo diffeomorphism,
and prove lens rigidity under a convex-foliation hypothesis. The unknown
there is a Riemannian metric. Our metric is fixed Euclidean, the unknowns
are periodic reflecting bodies and their translation group, and the
experiment includes internal localized impacts. Neither information set
is silently identified with the other.

Gurfinkel--Noakes--Stoyanov~\cite{GNS23a}, Theorem 1, prove uniqueness
for finite unions of strictly convex obstacles in a containing strictly
convex Riemannian domain under their curvature conditions and equivalence
up to tangency; general position implies the latter condition. Their
subsequent work~\cite{GNS23b}, Theorem 1, removes that tangency-equivalence
hypothesis in dimension at least three under the stated curvature
conditions. Its two-dimensional result is conditional on the admissible-path
property in Theorem 2. These are exterior travelling-time data, triples
of entry point, exit point and length on a containing boundary. We do not
claim their higher-dimensional result without qualifications in the plane,
nor treat our collision-localization apparatus as an exterior travel-time
sensor. Their results do not supply the particular periodic completion
defect or the finite launch accounting studied here; obstacle recovery
from rich scattering information itself is not a novelty claim.

The collision-cloud proof also has a direct stochastic-geometry comparison.
Schneider~\cite{Schneider} studies boundary-sampled random convex hulls in
the plane, including Hausdorff approximation. Prochno--Schuett--Sonnleitner--Werner
\cite{PSSW} give modern Hausdorff-distance results and record the smooth
boundary order $(\log n/n)^2$ in dimension two. We neither claim that
boundary-hull rate nor mollifier interpolation as new. Our use proves
that the requisite boundary coverage is generated by the physical
first-hit experiment without known labels, and combines it with
unknown-solid confidence, period recognition and refreshed stopping.
The finite bound here is deliberately conservative and is not an
asymptotic constant or minimax theorem.

\subsection{Qualification of the retained effective theorem}
To make the earlier numerical-prior convention unambiguous, its effective
recognition theorem uses the following directed inputs: centered support
strip width and upper bound $(\rho,M)$; local graph disk radius and upper
bound $(R,M_g)$; lower curvature $\kappa_-$; upper bridge length $R_+$
and body diameter $D_0$; lower rotation/reflection margin $\eta$, inter-type
shape gap $\sigma$, and physical-body separation $d_0$; and all chart,
time, active and clearance margins required for acquisition. Lower bounds
must be positive and valid, and upper bounds must actually dominate the
corresponding quantities. These analytic inputs belong to that retained
short-arc theorem, not to Definition~\ref{def:prior}.

In the retained value-probe cost statement, the required height error
\emph{tolerance is at most} $c\nu^3$. The ambiguous expression about
accuracy being ``no smaller'' is to be read with this explicit inequality.
The retained per-probe constant contains its fixed value-grid size $N$;
so do its epoch-work constants. Expressions such as $N=2^P$ describe
symbolic computability, not a physically expanded or queried vector.
The new descriptor has size \eqref{eq:fingerprint-size} instead, but
uses measured full-body geometry. Comparing the two sizes without that
sensor distinction would be misleading.

The old order-dependent moment regularization, analytic full-histogram
regularization and present cloud-assisted regularization remain different
experiments. The first uses finitely many types of exact moment germs;
the second measures densities and continues arcs; the third samples full
nearby bodies as well as the normalized endpoint law. No optimality
comparison follows by comparing powers of their different error parameters.

\subsection{Preserved manuscript and scope}
The complete reviewed v25 manuscript is supplied unchanged in
\texttt{retained/v25/}, with every earlier retained manuscript nested
inside it. It includes the analytic effective calibration, local-value
construction, coefficient inverses, count fibers, registered and
unregistered variants, and their original proofs. The original complete
smooth relative-law supplement is also preserved at \texttt{complete/}.
The source map identifies the active main documents; previous build
receipts remain historical and are not evidence of a current execution.
This primary article is self-contained for its new finite-smoothness
certification route. No retained mathematical argument is deleted.

Finally, equality of finite or formal count data, equality of all smooth
Taylor coefficients, and equality of exact analytic count germs are not
interchangeable. Neither the new collision cloud nor its completion
certificate proves a count-only analytic theorem. Nor is the result a
marked-length-spectrum rigidity theorem, a passive-trajectory discovery
theorem, a theorem without positive generic margins, or a claim of
journal acceptance. The new assertions concern the explicitly stated
launch experiment and its reference-free whole-table certificate.
